\scanpage{77}
\mspage{68}
% archive sheet number: 77

\[
\Sigma_M\longrightarrow\Sigma_L\qquad S\longrightarrow S\cap L
\]
mais est-on sûr que si $S$ est un support\inserted{(interne)}, $S\cap L$ en soit un ? C'est OK si on suppose Mag L 2 \deleted{\unsure{qui} \illegible{2} \unsure{dans le cas quasi-accessible}} \inserted{(\unsure{magmas Art. dirigibles|magmas Art. divisibles})}, car alors,
\[
\operatorname{cosupp}^{\circ}_M X=\operatorname{cosupp}^{\circ}_M(\operatorname{omb}^{\circ}X)
\]
% margin note written slanted
\margin{NB. On en conclut\\$\operatorname{supp}^{\circ}X=\operatorname{supp}^{\circ}(\operatorname{omb}^{\circ}X)$}
et plus généralement, pour une partie quelc. $A$ de $M$
\[
\operatorname{cosupp}^{\circ}_M(A)=\operatorname{cosupp}^{\circ}_M(\operatorname{omb}^{\circ}(A))
\]
où on pose
\[
\operatorname{omb}^{\circ}(A)=\bigcup_{X\in A}\operatorname{omb}^{\circ}(X)\subset L
\]
on en conclut
\[
\operatorname{cosupp}^{\circ}_M(A)\cap L=\operatorname{cosupp}^{\circ}_L(\operatorname{omb}(A)),
\]
\deleted{\unsure{donc}} ce qui prouve que si $S$ est un support\inserted{(interne)} dans $M$ i.e. de la forme $\operatorname{cosupp}^{\circ}_M(A)$, alors $S\cap L$ est un support (interne) dans $L$.

On trouve une application en sens inverse
\[
\begin{aligned}
\Sigma_L&\longrightarrow\Sigma_M\\
 A&\longrightarrow\operatorname{supp}^{\circ}_M(A)\underset{\uparrow}{\overset{?}{=}}\{X\in M\mid\operatorname{omb}^{\circ}X\subset A\unsure{)|\}}
\end{aligned}
\]
A-t-on bien l'égalité annoncée ? Sûrement, si
\[
\inserted{X\in M}\ \deleted{\unsure{si}}\quad\operatorname{omb}^{\circ}(X)\subset A\quad\Rightarrow\quad
\underset{\inserted{\substack{\Vert\\\operatorname{supp}^{\circ}(\operatorname{omb}^{\circ}(X))}}}{\deleted{\unsure{\operatorname{supp}}}\operatorname{supp}^{\circ}(X)}
\unsure{\subset|\in}\operatorname{supp}^{\circ}_M(A)
\]
d'où $X\in\deleted{\unsure{\operatorname{cosupp}}}\operatorname{supp}^{\circ}_M(A)$. Mais inversement, supposons
